\documentclass[
  12pt,			      % tamanho da fonte
  openright,		      % capítulos começam em pág ímpar (insere página vazia caso preciso)
  oneside,			  % para impressão em recto e verso. Oposto a oneside
  a4paper,			  % tamanho do papel. 
  chapter=TITLE,        % títulos de capítulos convertidos em letras maiúsculas
  section=TITLE,        % títulos de seções convertidos em letras maiúsculas
  english,			  % idioma adicional para hifenização
  french,				  % idioma adicional para hifenização
  spanish,			  % idioma adicional para hifenização
  brazil				  % o último idioma é o principal do documento
]{iftex}

% Importa as macros do projeto
%% Macros de dados do documento
\autor{Henrique Machado de Oliveira}
\titulo{Redes Neurais Convolucionais para o Controle de Qualidade de Rochas Ornamentais: Da Classificação à Segmentação de Anomalias utilizando pipeline com abordagem Teacher-Student}

%%%Sistema Hierárquico de Inteligência Artificial para o Controle de Qualidade de Rochas Ornamentais: Da Classificação à Segmentação de Irregularidades
%%%Análise Comparativa entre Modelos de Segmentação Generalistas e Arquiteturas em Cascata para Identificação de Defeitos em Rochas Ornamentais
%%%Otimização da Segmentação de Superfícies Pétreas através de Classificação Prévia e Redes Especialistas
%%%Abordagem Hierárquica em Deep Learning: Impacto da Especialização de Modelos na Precisão da Segmentação de Rochas Ornamentais

\instituicao{Instituto Federal do Espírito Santo}
\curso{Bacharelado em Sistemas de Informação}
\local{Cachoeiro de Itapemirim}
\data{2026}
\tipotrabalho{Trabalho de Graduação}
\preambulo{Trabalho de Conclusão de Curso apresentado à Coordenadoria do Curso de Sistemas de Informação do Instituto Federal de Educação, Ciência e Tecnologia do Espírito Santo, como requisito parcial para a obtenção do título de Bacharel em Sistemas de Informação.}

\orientador[Orientador][Prof. Dr.]{Rafael}[Silva Guimarães]{Instituto Federal do Espírito Santo}
%%%\coorientador[Coorientador][Prof. Dr.]{Everson}[Scherrer Borges]{Instituto Federal do Espírito Santo}
%%%\coorientador[Coorientador][Prof. Dr.]{Igor}[Henrique Beloti Pizetta]{Instituto Federal do Espírito Santo}

\examinadori{Profa. Dra. Fulana de Tal}{Instituto Federal do Espírito Santo}{Examinadora}
\examinadorii{Prof. Dr. Cicrano de Tal}{Instituto Federal do Espírito Santo}{Examinador}

% TODO: data real da defesa, a confirmar com o orientador (ver docs/pendencias.md).
\approvaldate{15}{Outubro}{2026}

% Palavras-chaves
\palavraschave{Visão Computacional}[Segmentação Semântica][Rochas Ornamentais][Modelos de Fundação][Pseudo-rotulagem]
\keywords{Computer Vision}[Semantic Segmentation][Dimension Stone][Foundation Models][Pseudo-labeling]

% Ficha catalográfica
\fichabiblioteca{
    Dados Internacionais de Catalogação-na-Publicação (CIP) \\
    (Biblioteca Nilo Peçanha do Instituto Federal do Espírito Santo)
}
\fichaautor{Elaborada por XXXXXXXXXXXXXXXXXXXXX – CRB-X/ES - XXX}
\fichatags{
    1. Visão computacional.  
    2. Aprendizado profundo.  
    3. Segmentação de imagens.  
    4. Rochas ornamentais – Controle de qualidade.  
    5. Inteligência artificial.  
    I. Guimarães, Rafael Silva.  
    II. Instituto Federal do Espírito Santo.  
    III. Título.
}

\cutter{X999y}
\cdd{000.00}
\selectlanguage{brazil}

% Arial (para usar times-new-roman, comente as três linhas abaixo)
\usepackage{helvet}
\renewcommand{\familydefault}{\sfdefault}
\renewcommand{\ABNTEXchapterfont}{\bfseries}
% Comandos para revisar o texto
\usepackage{easyReview} 


%% Elementos opcionais
\editardedicatoria{A dedicatória é um elemento opcional. Contém o oferecimento do trabalho adeterminada pessoa ou a pessoas (ASSOCIAÇÃO BRASILEIRA DE NORMAS TÉCNICAS, 2011b).}
\editaragradecimentos{A um elemento opcional. Localiza-se após a folha de aprovação e deve ser dirigido àqueles que realmente contribuíram, de maneira relevante, para a elaboração do trabalho. Deve-se utilizar uma linguagem simples (ASSOCIAÇÃO BRASILEIRA DE NORMAS TÉCNICAS, 2011b).}
\editarepigrafe{\textbf{Epigrafe:} É um elemento opcional. É uma citação relacionada ao assunto do trabalho desenvolvido, seguida da indicação de autoria (ASSOCIAÇÃO BRASILEIRA DE NORMAS TÉCNICAS, 2011b).   Deve-se seguir as regras do uso da citação NBR 10.520/2002.}


%% Listas [Siglas e Simbolos]
\editarlistasiglas{\begin{siglas}
    \item[ARIA] Análise e Reconhecimento Inteligente de Anomalias
    \item[CLIP] \textit{Contrastive Language--Image Pre-training}
    \item[CNN] Rede Neural Convolucional (\textit{Convolutional Neural Network})
    \item[GPU] Unidade de Processamento Gráfico (\textit{Graphics Processing Unit})
    \item[Ifes] Instituto Federal do Espírito Santo
    \item[IoU] \textit{Intersection over Union}
    \item[mAP] \textit{mean Average Precision}
    \item[NMS] \textit{Non-Maximum Suppression}
    \item[SAM] \textit{Segment Anything Model}
    \item[YOLO] \textit{You Only Look Once}
\end{siglas}
}
\editarlistasimbolos{\begin{simbolos}
    \item[$\lambda$] Letra grega labda
\end{simbolos}}

% Novos comandos
\newcommand{\ie}{\textit{i.e. }}
\newcommand{\eg}{\textit{e.g. }}

\begin{document}

  % --------------------------------------------------- %
  % CAPAS                                               %
  % --------------------------------------------------- %
  % CAPA E FOLHA DE ROSTO
  \capas
  \imprimircapa
  \imprimirfolhaderosto
  \imprimirfichacatalografica
  % Quando a biblioteca devolver a ficha oficial em PDF, troque a linha acima por:
  % \includepdf[pages=-]{pre_textuais/ficha.pdf}

  % --------------------------------------------------- %
  % ELEMENTOS PRÉ-TEXTUAIS                              %
  % --------------------------------------------------- %
  % APROVAÇÃO, DEDICATÓRIA, AGRADECIMENTOS, EPÍGRAFE
  \pretextual
  % TODO: os examinadores em macros.tex ainda são "Fulana/Cicrano de Tal".
  % Descomente a aprovação quando a banca estiver definida.
  %%%\imprimiraprovacao
  % \includepdf[pages=-]{pre_textuais/aprovacao.pdf}
  % TODO: dedicatória, agradecimentos e epígrafe ainda são texto do template.
  %%%\imprimirdedicatoria
  %%%\imprimiragradecimentos
  %%%\imprimirepigrafe
  
  % RESUMO - PT/EN
  % RESUMO - PT
\begin{resumo}
  \vspace{-15pt}

  A inspeção de anomalias superficiais em chapas de rochas ornamentais é feita visualmente por operadores, e o critério que separa uma feição natural de um defeito comercial não está escrito em lugar nenhum: ele varia com a rocha, o cliente, o lote e o inspetor. Este trabalho parte da premissa de que esse critério é arbitrário por natureza e de que a contribuição possível não é eliminá-lo, mas torná-lo explícito, parametrizado e auditável. Propõe-se o ARIA (Análise e Reconhecimento Inteligente de Anomalias), um \textit{pipeline} hierárquico de visão computacional em que um classificador identifica a litologia da chapa, um modelo de fundação guiado por texto atua como Professor gerando polígonos de anomalia sem anotação humana, e uma rede de segmentação em tempo real atua como Aluno na linha de produção. O critério de marcação é materializado em um conjunto de sondas textuais e limiares de confiança por litologia, gravado em arquivo versionável. A calibração segue um protocolo de quatro imagens por litologia --- uma para descobrir quais sondas respondem e três para submeter um único limiar a casos sutil, típico e forte --- e o limiar é produzido por uma regra idêntica para todas as litologias. Demonstra-se que o limiar de confiança do Professor é um filtro aplicado após a inferência, o que torna a varredura de limiares fora de linha exatamente equivalente a reexecutar o modelo em cada valor; a equivalência foi verificada experimentalmente. Dois experimentos avaliam a proposta: o primeiro compara a mesma regra aplicada por litologia e ao conjunto reunido; o segundo compara modelos Aluno especialistas e generalista, estratificados por volume de dados disponível por litologia. A avaliação usa um conjunto de referência anotado antes de qualquer inferência sobre as mesmas imagens.
  % TODO: acrescentar uma a duas frases de RESULTADOS e CONCLUSÃO quando os
  % experimentos rodarem. A ABNT pede resumo informativo (objetivo, método,
  % resultados, conclusão) -- hoje só os dois primeiros existem de fato.

  Palavras-chave: \palavraschaveemlinha
\end{resumo}


% RESUMO - EN
\begin{resumo}[Abstract]
  \vspace{-15pt}

  \begin{otherlanguage*}{english}
  Surface anomaly inspection of polished dimension stone slabs is performed visually by human operators, and the criterion separating a natural feature from a commercial defect is nowhere written down: it varies with the stone, the customer, the batch and the inspector. This work takes that criterion to be arbitrary by nature and argues that the available contribution is not to eliminate it, but to make it explicit, parameterised and auditable. We propose ARIA, a hierarchical computer vision pipeline in which a classifier identifies the slab lithology, a text-prompted foundation model acts as a Teacher producing anomaly polygons without human annotation, and a real-time segmentation network acts as a Student on the production line. The marking criterion is materialised as a set of textual probes and confidence thresholds per lithology, stored in a versionable file. Calibration follows a four-image protocol per lithology --- one image to discover which probes respond, and three to subject a single threshold to a subtle, a typical and a strong case --- with the threshold produced by a rule that is identical across all lithologies. We show that the Teacher's confidence threshold is a filter applied after inference, which makes an offline threshold sweep exactly equivalent to re-running the model at each value; the equivalence was verified experimentally. Two experiments evaluate the proposal: the first compares the same rule applied per lithology against the pooled material; the second compares specialist and generalist Student models, stratified by the data volume available per lithology. Evaluation uses a reference set annotated before any inference was run on those same images.
  % TODO: mirror the RESULTS/CONCLUSION sentences added to the Portuguese abstract.

  Keywords: \inlinekeywords
\end{otherlanguage*}
\end{resumo}

  
  % LISTAS E SUMÁRIO
  \imprimirlistafiguras
  \imprimirlistatabelas
  \imprimirlistaquadros
  \imprimirlistasiglas
  % Sem lista de símbolos: o trabalho não usa notação simbólica própria.
  %%%\imprimirlistasimbolos
  \imprimirsumario

  % --------------------------------------------------- %
  % ELEMENTOS TEXTUAIS                                  %
  % --------------------------------------------------- %
  \textual
  \graphicspath{{figuras/}}   
  \chapter[INTRODUÇÃO]{INTRODUÇÃO}

\section{Contextualização}

O Espírito Santo responde pela maior parte da produção e da exportação brasileira de rochas ornamentais, atividade que sustenta uma cadeia industrial extensa, da lavra ao beneficiamento \cite{stone_economy_bahiense2023}. Ao fim dessa cadeia, cada chapa polida passa por uma etapa de inspeção visual em que um operador decide se a peça segue para o cliente, é reclassificada ou é descartada. Essa decisão determina o preço da chapa e, agregada, uma fração relevante da margem da empresa \cite{stone_inspection_vidal2013}.

A inspeção é feita a olho nu, sob iluminação variável e em ritmo de produção. O operador avalia a superfície da chapa procurando trincas, manchas, oxidações e descontinuidades que comprometam o uso comercial da peça. O problema não está na acuidade visual de quem inspeciona, mas em algo anterior: \textbf{não existe um critério escrito que separe uma feição natural da rocha de um defeito comercial}. A fronteira entre as duas categorias depende da litologia, do cliente, do lote e do acabamento contratado --- e, na prática, depende do inspetor \cite{stone_subjectivity_castro2018}.

Essa ausência tem uma consequência que costuma passar despercebida: o critério existe, mas vive implícito na experiência de cada operador. Ele não pode ser consultado, discutido, comparado entre turnos nem auditado por um cliente que conteste uma reclassificação. Duas chapas equivalentes podem receber destinos diferentes sem que haja registro de qual regra produziu cada decisão.

A visão computacional é uma resposta natural a esse tipo de problema, e o setor já registra tentativas de automação da inspeção de superfícies pétreas \cite{stone_tereso2020, stone_lopez2010}. A dificuldade é que as abordagens supervisionadas convencionais exigem grandes volumes de imagens anotadas manualmente, e anotar anomalias sutis em chapas de rocha é caro, lento e --- pelo mesmo motivo descrito acima --- subjetivo. Anotar milhares de imagens apenas transfere a arbitrariedade do chão de fábrica para o conjunto de treinamento.

\section{Problema de Pesquisa}

O problema que este trabalho ataca não é a dificuldade técnica de detectar anomalias, e sim a \textbf{arbitrariedade não registrada do critério de marcação}. Não existe uma resposta absolutamente correta a ser descoberta: o que é defeito em uma litologia pode ser exatamente o padrão estético valorizado em outra. Admitir isso muda o que se pode legitimamente prometer. Um sistema automatizado não elimina a arbitrariedade da fronteira; o que ele pode fazer é retirá-la da cabeça do inspetor e colocá-la em um parâmetro inspecionável.

Formula-se, assim, a seguinte pergunta de pesquisa: \textit{é possível materializar o critério de marcação de anomalias em chapas de rochas ornamentais como um conjunto de parâmetros explícito, versionável e específico por litologia --- e esse critério parametrizado produz marcações mais alinhadas a uma referência anotada do que um critério único aplicado a todas as rochas?}

\section{Hipóteses}

Duas hipóteses orientam o trabalho:

\begin{description}
  \item[H1] Um critério de marcação parametrizado por litologia produz segmentações mais alinhadas ao critério de referência anotado do que um critério único e global.
  \item[H2] Modelos de segmentação especialistas --- um por litologia --- alinham-se melhor à referência do que um único modelo generalista, e esse ganho é função do volume de dados disponível por litologia.
\end{description}

A segunda metade de H2 é uma pergunta de pesquisa por direito próprio: \textit{quantas imagens uma litologia precisa ter para que a especialização compense o custo de manter um modelo dedicado?} O conjunto de dados utilizado apresenta desbalanceamento acentuado entre litologias, o que permite tratar o volume de dados como variável do experimento em vez de limitação a ser contornada.

\section{Objetivos}

\subsection{Objetivo Geral}

Desenvolver e avaliar o ARIA (Análise e Reconhecimento Inteligente de Anomalias), um \textit{pipeline} hierárquico de visão computacional que materializa o critério de marcação de anomalias superficiais em chapas de rochas ornamentais como um conjunto explícito e versionável de parâmetros por litologia, e verificar se essa parametrização produz marcações mais alinhadas a uma referência anotada do que um critério único global.

\subsection{Objetivos Específicos}

\begin{enumerate}
  \item Selecionar, caracterizar e pré-processar um conjunto de dados público de imagens de chapas polidas de rochas ornamentais, descrevendo sua composição, o desbalanceamento entre litologias e o particionamento adotado.
  \item Definir um protocolo reprodutível de calibração que determine, para cada litologia, quais sondas textuais entram no conjunto e qual limiar de confiança cada uma recebe, de modo que a regra que produz o limiar seja a mesma para todas as litologias.
  \item Implementar o estágio Professor com um modelo de fundação guiado por texto, capaz de gerar polígonos de anomalia sem anotação humana prévia.
  \item Construir um conjunto de referência anotado independentemente das saídas do modelo, para servir de gabarito quantitativo da avaliação.
  \item Avaliar experimentalmente a hipótese H1, comparando o critério parametrizado por litologia ao mesmo critério derivado do conjunto reunido de litologias.
  \item Especificar o desenho experimental que avalia a hipótese H2, treinando modelos Aluno especialistas e generalista estratificados por faixa de volume de dados.
\end{enumerate}

\section{Justificativa}

Do ponto de vista industrial, tornar explícito o critério de marcação tem valor independente da acurácia alcançada. Um arquivo de parâmetros pode ser versionado, comparado entre unidades produtivas, negociado com um cliente e auditado quando uma classificação é contestada --- coisas que um critério que vive na experiência de um operador não permite. A automação da marcação é consequência; a contribuição central é a explicitação.

Do ponto de vista acadêmico, o trabalho se situa no encontro de duas linhas ativas. A primeira é o uso de modelos de fundação como geradores de rótulos para domínios sem \textit{ground truth}, o que dispensa a anotação manual em larga escala. A segunda é o aprendizado a partir de rótulos ruidosos, campo que estuda justamente o que acontece quando a supervisão vem de um modelo e não de um humano. O trabalho contribui com um elemento pouco explorado nessa interseção: o \textbf{condicionamento do \textit{prompt} por classe do domínio} --- aqui, por litologia --- como forma de controlar a qualidade do rótulo gerado, e o registro desse condicionamento como artefato auditável.

Por fim, o uso de um conjunto de dados público \cite{dataset_araujo2023} garante que todo o protocolo possa ser reproduzido por terceiros sem acordo com nenhuma empresa, o que é incomum em trabalhos de inspeção industrial.

\section{Organização do Trabalho}

O restante deste documento está organizado da seguinte forma. O Capítulo 2 apresenta a fundamentação teórica, cobrindo a inspeção de qualidade no setor de rochas ornamentais, redes neurais convolucionais aplicadas à classificação de litologias, modelos de fundação com segmentação guiada por texto, a arquitetura Professor--Aluno com pseudo-rotulagem, segmentação em tempo real, e os trabalhos relacionados contra os quais a contribuição se posiciona. O Capítulo 3 detalha a metodologia: o conjunto de dados, a arquitetura do \textit{pipeline}, o protocolo de calibração, a regra que produz o limiar, a estratégia de varredura de limiares fora de linha, a referência de avaliação e o desenho dos dois experimentos. O Capítulo 4 apresenta os resultados obtidos. O Capítulo 5 traz as considerações finais, as limitações reconhecidas e os desdobramentos deixados como trabalho futuro.

  \chapter[FUNDAMENTAÇÃO TEÓRICA]{FUNDAMENTAÇÃO TEÓRICA}

Este capítulo reúne os fundamentos que sustentam o ARIA. Parte da natureza do problema de inspeção no setor de rochas ornamentais, percorre as arquiteturas de visão computacional empregadas em cada estágio do \textit{pipeline}, e detém-se no ponto que mais condiciona o desenho do trabalho: o que a literatura sabe sobre treinar uma rede a partir de rótulos que não foram produzidos por um humano.

\section{Inspeção de Qualidade em Rochas Ornamentais}

A inspeção de chapas polidas é uma etapa de classificação comercial, não de detecção de falhas no sentido estrito. Uma mesma feição --- um veio mineral marcante, uma variação de tonalidade --- pode ser o atributo que valoriza uma chapa e o que desqualifica outra, conforme a litologia e o mercado de destino \cite{stone_inspection_vidal2013}. Por isso a inspeção depende de julgamento acumulado, e não da aplicação de uma norma objetiva.

Trabalhos que documentam a prática do setor registram a variabilidade dessa decisão entre inspetores e entre turnos \cite{stone_subjectivity_castro2018}. Tentativas de automação por visão computacional existem há tempo e concentram-se na classificação da rocha ou na detecção de descontinuidades específicas \cite{stone_lopez2010, stone_tereso2020}. O contexto mais amplo de digitalização da manufatura, associado à Indústria 4.0, fornece a motivação econômica para essa automação \cite{industry40_lee2014}.

Em outros domínios industriais, a inspeção automática de superfícies é assunto consolidado, com aplicações em soldagem \cite{avi_welding_tsuzuki2022}, concreto \cite{avi_concrete_waqar2024}, têxteis \cite{avi_textile_sikka2024} e alimentos \cite{avi_food_chen2021}, e revisões recentes consolidam as abordagens de detecção de anomalias visuais \cite{avi_anomaly_li2025}. A diferença do caso das rochas ornamentais é que nesses domínios o defeito costuma ter definição física estável --- uma trinca em uma solda é uma trinca --- enquanto aqui a categoria "defeito" é, em parte, uma convenção comercial.

\section{Redes Neurais Convolucionais e Classificação de Litologias}

Redes neurais convolucionais aprendem hierarquias de representações diretamente dos pixels, dispensando a engenharia manual de descritores \cite{deeplearning_lecun2015}. A arquitetura Xception substitui as convoluções tradicionais por convoluções separáveis em profundidade, o que reduz o número de parâmetros sem perda de capacidade representacional \cite{xception_chollet2017}.

A identificação automática de tipos de rocha a partir de imagens é tarefa estabelecida na literatura geológica computacional \cite{rocktyping_baraboshkin2020}. No domínio específico das rochas ornamentais brasileiras, trabalhos anteriores do mesmo grupo de pesquisa aplicaram redes convolucionais ao reconhecimento de litologias em chapas polidas \cite{deepstone_araujo2022}, resultando no classificador DeepStoneAI \cite{deepstoneai_moreira2025}, que é o modelo adotado como primeiro estágio deste trabalho.

\section{Modelos de Fundação e Segmentação Guiada por Texto}

Modelos de fundação são redes treinadas em conjuntos de dados de escala muito superior à de um domínio específico, capazes de executar tarefas para as quais não foram explicitamente treinadas --- a chamada capacidade \textit{zero-shot}. Em segmentação, o marco dessa mudança foi o \textit{Segment Anything Model}, que introduziu a segmentação condicionada por \textit{prompt}: em vez de aprender um conjunto fixo de classes, o modelo recebe uma indicação --- um ponto, uma caixa, um texto --- e devolve a máscara correspondente \cite{sam_kirillov2023}. Este trabalho emprega a terceira geração dessa família, que aceita \textit{prompt} textual de forma nativa \cite{TODO-sam3}.

A ponte entre texto e imagem é feita por representações multimodais alinhadas, treinadas de modo contrastivo sobre pares imagem--legenda \cite{clip_radford2021}. É essa representação compartilhada que permite que uma palavra selecione uma região da imagem sem nenhum treinamento supervisionado no domínio de destino.

Um achado relevante do trabalho original sobre essas representações é que \textit{prompts} contextuais superam consistentemente palavras isoladas em tarefas \textit{zero-shot}: o modelo responde melhor a \textit{a photo of a \{classe\}} do que a \textit{\{classe\}} \cite{clip_radford2021}. A observação é diretamente aplicável ao domínio deste trabalho, em que termos como \textit{vein} são lexicalmente ambíguos, e é registrada aqui como atribuição correta do achado --- não como hipótese original desta pesquisa.

\section{Professor--Aluno e Pseudo-rotulagem}

A transferência de conhecimento de um modelo grande para um modelo menor é ideia consolidada desde os trabalhos sobre destilação \cite{distill_hinton2015}. A pseudo-rotulagem é a variante em que o modelo grande não transfere distribuições de saída, mas produz \textbf{rótulos} para dados não anotados, que passam a servir de supervisão para o modelo menor \cite{pseudolabel_lee2013}.

No arranjo adotado neste trabalho, o modelo de fundação atua como Professor, processando as imagens fora da linha de produção e gerando polígonos de anomalia; uma rede de segmentação compacta atua como Aluno, treinada sobre esses polígonos e responsável pela inferência em produção. A vantagem prática é econômica: o custo de inferência do modelo de fundação é pago uma única vez, na geração do conjunto de treinamento, e não a cada chapa inspecionada.

\section{Aprendizado com Rótulos Ruidosos e a Escolha do Limiar}

O rótulo produzido por um Professor é, por construção, ruidoso. Essa constatação define o problema central da calibração: \textbf{onde cortar o escore de confiança do Professor}.

A literatura de detecção semissupervisionada converge em uma direção. \citeonline{softteacher_xu2021} reportam melhor desempenho com limiares de confiança altos, observando o compromisso esperado --- limiar mais alto aumenta a precisão do pseudo-rótulo e reduz o \textit{recall}. \citeonline{unbiasedteacher_liu2021} chegam a conclusão análoga ao tratar o viés de confirmação introduzido por pseudo-rótulos incorretos. A razão de fundo é conhecida: redes de alta capacidade \textbf{memorizam} rótulos errados quando expostas a eles repetidamente \cite{memorization_arpit2017}, de modo que um rótulo falso-positivo não é apenas ruído aleatório, mas um sinal que a rede efetivamente aprende.

Duas ressalvas importam para a transposição desse consenso ao presente trabalho. A primeira é que os escores relatados naqueles trabalhos vêm de detectores supervisionados com saídas calibradas, e não são comparáveis em magnitude aos escores de um modelo de fundação guiado por texto --- razão pela qual este trabalho transporta a \textit{direção} da recomendação (na dúvida, cortar mais apertado) e não o valor numérico. A segunda é que o compromisso precisão--\textit{recall} do limiar não tem solução ótima única quando as imagens de uma mesma litologia variam em intensidade das feições; daí o protocolo de calibração descrito no Capítulo 3 submeter um mesmo limiar a três casos distintos em vez de otimizá-lo em uma imagem.

Para converter uma curva de comportamento em um ponto de operação, este trabalho adota a detecção de joelho proposta por \citeonline{kneedle_satopaa2011}, que define o joelho como o ponto de máxima distância à corda que liga os extremos da curva normalizada.

\section{Segmentação em Tempo Real}

A segmentação semântica e de instâncias evoluiu de arquiteturas encoder--decoder \cite{unet_ronneberger2015} e de detectores em dois estágios estendidos para máscaras \cite{maskrcnn_he2017} até famílias de estágio único otimizadas para latência \cite{segsurvey_minaee2021}. A família YOLO consolidou a abordagem de predição em uma única passagem pela rede \cite{yolo_redmon2016}, e suas versões recentes incorporam cabeças de segmentação mantendo taxas de quadros compatíveis com inspeção em linha \cite{yolo11_khanam2024}. É essa característica --- inferência rápida em \textit{hardware} de custo moderado --- que justifica seu papel como Aluno no \textit{pipeline}.

\section{Métricas de Avaliação}

A sobreposição entre uma região predita e uma região de referência é medida pela \textit{Intersection over Union} (IoU), definida na Equação \ref{eq:iou} \cite{iou_mahmood2021}:

\begin{equation}
  \mathrm{IoU}(A, B) = \frac{|A \cap B|}{|A \cup B|}
  \label{eq:iou}
\end{equation}

\noindent onde $A$ é a região predita e $B$ a região de referência. A precisão média (AP) agrega precisão e \textit{recall} ao longo dos limiares de decisão, e a \textit{mean Average Precision} (mAP) é sua média sobre as classes \cite{voc_everingham2010}. Como este trabalho adota rotulagem binária --- justificada na Seção~\ref{sec:sondas} ---, há uma única classe de interesse e a mAP degenera na própria AP; o texto usa AP para evitar a sugestão falsa de uma média sobre múltiplas classes.

\section{Trabalhos Relacionados e Posicionamento}
\label{sec:relacionados}

O trabalho mais próximo do aqui proposto emprega máscaras de um modelo de fundação como pseudo-rótulos para treinar um segmentador de defeitos industriais, tratando explicitamente o ruído dessas máscaras \cite{boxes2pixels_lendering2026}. A coincidência de estrutura é considerável: mesmo arranjo Professor--Aluno, mesma origem de rótulo, mesmo domínio geral de inspeção industrial.

Convém, portanto, enunciar o posicionamento deste trabalho sem recorrer à alegação de lacuna. A distinção não está em usar um modelo de fundação como Professor --- isso já foi feito ---, e sim em dois pontos: \textbf{(i)} o \textit{prompt} do Professor é condicionado à classe do domínio, de modo que cada litologia recebe seu próprio conjunto de sondas e seus próprios limiares, em vez de uma configuração única para todo o conjunto de dados; e \textbf{(ii)} essa configuração é tratada como o artefato central do trabalho --- a materialização auditável de um critério que, no processo manual, não existe por escrito. O domínio de rochas ornamentais, em que a fronteira entre feição natural e defeito é convenção comercial e não propriedade física, é o que torna o ponto (ii) mais do que uma conveniência de engenharia.

  \chapter[METODOLOGIA]{METODOLOGIA}

Este capítulo descreve o conjunto de dados, a arquitetura do \textit{pipeline}, o protocolo de calibração, a regra que produz os limiares, a referência adotada para avaliação e o desenho dos dois experimentos.

\section{Caracterização da Pesquisa}

O trabalho classifica-se como pesquisa aplicada, de natureza experimental e abordagem predominantemente quantitativa. É aplicada porque endereça um problema concreto da indústria de rochas ornamentais. É experimental porque manipula variáveis controladas --- o escopo sobre o qual a regra de calibração é aplicada e a quantidade de modelos treinados --- e observa seu efeito sobre métricas de sobreposição com uma referência anotada.

\section{Conjunto de Dados}

O conjunto de imagens é público, obtido na plataforma Kaggle \cite{dataset_araujo2023}. O autor deste trabalho \textbf{não constituiu} o banco: selecionou, caracterizou e pré-processou um conjunto existente. Essa escolha tem uma consequência metodológica favorável --- todo o protocolo pode ser reproduzido por terceiros sem acordo com nenhuma empresa.

O conjunto reúne 34.630 imagens de chapas polidas, distribuídas em 45 litologias. O particionamento é aleatório estratificado nas proporções de 70\%, 15\% e 15\%, aplicado a todas as classes, resultando em 24.263 imagens de treinamento, 5.214 de validação e 5.153 de teste.

A distribuição entre litologias é fortemente desbalanceada, variando de 106 a 4.588 imagens. Em vez de equalizar artificialmente essa distribuição, o trabalho a converte em variável do experimento, agrupando as litologias em quatro faixas de volume, conforme o Quadro \ref{quad:faixas}.

\begin{quadro}[ht]
  \caption{Faixas de volume de dados e respectivo número de litologias.}
  \label{quad:faixas}
  \centering
  \begin{tabularx}{\textwidth}{clcX}
    \hline
    \bfseries Faixa & \bfseries Critério & \bfseries Litologias & \bfseries Observação \\
    \hline
    A & $\geq$ 1.000 imagens & 11 & Faixa de execução prioritária \\
    B & 500 a 999 imagens & 6 & \\
    C & 200 a 499 imagens & 14 & \\
    D & $<$ 200 imagens & 14 & Inclui as litologias de textura atípica \\
    \hline
  \end{tabularx}
  \legend{Fonte: elaborado pelo autor.}
\end{quadro}

A ordem de execução dos experimentos segue as faixas, da maior para a menor, e cada faixa constitui um marco entregável: ela é executada, avaliada e redigida antes que a seguinte comece. Em qualquer ponto de interrupção existe, assim, um resultado completo e reportável, e as faixas não alcançadas são declaradas como trabalho futuro.

\section{Arquitetura do \textit{Pipeline}}

O ARIA organiza-se em três estágios, descritos no Quadro \ref{quad:estagios}. O primeiro identifica a litologia da chapa e roteia; o segundo gera os polígonos de anomalia fora da linha de produção; o terceiro executa a inferência rápida em produção.

\begin{quadro}[ht]
  \caption{Estágios do \textit{pipeline} ARIA.}
  \label{quad:estagios}
  \centering
  \begin{tabularx}{\textwidth}{llX}
    \hline
    \bfseries Estágio & \bfseries Modelo & \bfseries Papel \\
    \hline
    1 --- Classificador & Xception (DeepStoneAI) & Identifica a litologia e seleciona a configuração de sondas e o modelo Aluno correspondente \\
    2 --- Professor & Modelo de fundação guiado por texto & Gera polígonos de anomalia sem anotação humana, processando o conjunto fora de linha \\
    3 --- Aluno & YOLO11-seg & Executa a inferência em produção, treinado sobre os polígonos do Professor \\
    \hline
  \end{tabularx}
  \legend{Fonte: elaborado pelo autor.}
\end{quadro}

A separação entre Professor e Aluno é o que torna o arranjo viável industrialmente: o custo de inferência do modelo de fundação é pago uma única vez, na geração do conjunto de treinamento, e não a cada chapa inspecionada.

% TODO: inserir a figura do pipeline quando o arquivo estiver em figuras/.
% O bloco abaixo está comentado de propósito: \includegraphics de um arquivo
% inexistente interrompe a compilação.
% \begin{figure}[ht]
%   \caption{Fluxo de processamento do \textit{pipeline} ARIA.}
%   \label{fig:arquitetura}
%   \centering
%   \includegraphics[width=0.9\textwidth]{arquitetura_aria.png}
%   \legend{Fonte: elaborado pelo autor.}
% \end{figure}

\section{Sondas de Recall, não Rótulos Semânticos}
\label{sec:sondas}

O Professor é guiado por um conjunto de palavras em inglês, aqui denominadas \textbf{sondas}. O Quadro \ref{quad:sondas} lista o conjunto de trabalho e a assinatura visual que cada sonda costuma ativar.

\begin{quadro}[ht]
  \caption{Sondas textuais, identificadores e assinatura visual associada.}
  \label{quad:sondas}
  \centering
  \begin{tabularx}{\textwidth}{llX}
    \hline
    \bfseries Sonda & \bfseries \textit{id} & \bfseries Assinatura visual que costuma ativar \\
    \hline
    \textit{vein}         & 0 & Estrias e faixas contrastantes \\
    \textit{crack}        & 1 & Descontinuidades lineares finas e escuras \\
    \textit{Stain}        & 2 & Variação de cor em mancha difusa \\
    \textit{Dark patches} & 3 & Regiões escuras sobre fundo claro \\
    \textit{light spot}   & 4 & Regiões claras sobre fundo escuro \\
    \textit{scratch}      & 5 & Riscos superficiais finos \\
    \hline
  \end{tabularx}
  \legend{Fonte: elaborado pelo autor.}
\end{quadro}

É necessário ser explícito quanto ao estatuto dessas palavras. Elas \textbf{não são afirmações sobre a natureza do que foi encontrado}: são chaves lexicais escolhidas por fazerem o modelo responder a certas assinaturas visuais. Afirmar que a região marcada pela sonda \textit{crack} \textit{é} uma fissura exigiria ter validado que o modelo distingue fissura de veio --- validação que este trabalho não realiza e, portanto, não pode alegar.

Dessa decisão decorre a \textbf{rotulagem binária}: todas as regiões marcadas recebem o mesmo identificador de classe no treinamento do Aluno. Os identificadores por sonda listados no Quadro \ref{quad:sondas} são preservados nos arquivos de anotação, para não perder o registro de qual sonda disparou cada região, mas são colapsados em uma única classe antes do treino. O objetivo do conjunto de sondas é maximizar a \textbf{cobertura} de regiões anômalas, não classificá-las. A rotulagem multiclasse fica registrada como trabalho futuro.

\section{Protocolo de Calibração}

Calibrar uma litologia significa responder a duas perguntas distintas: \textit{quais} sondas entram no conjunto, e com \textit{qual} limiar de confiança cada uma opera. A imagem ideal para responder a cada uma é oposta, e confundi-las é a origem de um viés concreto.

Para \textbf{descobrir} quais sondas respondem, a imagem deve conter o máximo de fenômenos --- quanto mais rica, melhor. Para \textbf{escolher} o limiar, a imagem deve representar a variação típica da litologia. Usar a chapa mais rica para as duas tarefas produz um efeito de seleção por visibilidade: as feições foram escolhidas justamente por serem visíveis, e o limiar resultante sai alto demais, perdendo o defeito sutil das demais chapas. Usar a chapa mais sutil produz o erro oposto --- limiar baixo e enxurrada de marcações na chapa comum.

O protocolo adotado separa as duas tarefas em quatro imagens por litologia, descritas no Quadro \ref{quad:vagas}.

\begin{quadro}[ht]
  \caption{As quatro imagens de calibração por litologia e o papel de cada uma.}
  \label{quad:vagas}
  \centering
  \begin{tabularx}{\textwidth}{llX}
    \hline
    \bfseries Imagem & \bfseries Decide & \bfseries Critério de seleção \\
    \hline
    \textit{descoberta}    & Quais sondas entram & A chapa com o maior número de tipos distintos de feição \\
    \textit{limiar sutil}  & O limiar & Feições fracas, de baixo contraste \\
    \textit{limiar típica} & O limiar & A chapa mais comum da litologia \\
    \textit{limiar forte}  & O limiar & Feições marcantes, de alto contraste \\
    \hline
  \end{tabularx}
  \legend{Fonte: elaborado pelo autor.}
\end{quadro}

Três observações completam o protocolo.

Primeiro, as três imagens de limiar \textbf{não produzem três limiares a serem promediados}. Trata-se de um único limiar submetido a três casos. A média de três ótimos não é um ótimo: o caso sutil pediria um valor baixo, o caso forte um valor alto, e a média seria um número nunca verificado em imagem alguma, possivelmente ruim nas duas pontas. A pergunta deixa de ser "qual limiar acerta esta imagem", que não tem resposta, e passa a ser "qual limiar menos erra nas três", que tem.

Segundo, a seleção das quatro imagens é \textbf{manual}. Usar o próprio Professor para escolher as imagens de calibração seria raciocínio circular.

Terceiro, as quatro imagens saem exclusivamente da partição de \textbf{treinamento}. O conjunto de referência da avaliação sai da partição de teste, e calibrar sobre uma imagem de teste equivaleria a ajustar o limiar sobre a própria prova. Como o particionamento é aleatório estratificado, as partições são amostras da mesma distribuição, de modo que a restrição não custa cobertura de fenômenos.

A origem de cada imagem selecionada --- partição e arquivo --- é registrada em um arquivo de metadados por litologia, compondo o material de reprodutibilidade do trabalho.

\section{A Regra que Produz o Limiar}
\label{sec:regra}

A escolha do limiar precisa ser governada por uma regra escrita. Selecionar o valor "no olho", ainda que por um observador experiente, reproduziria exatamente a prática que o Capítulo 1 aponta como problema --- com a diferença de que o observador passaria a ser o autor, em vez do inspetor.

Estabelece-se, portanto, a seguinte distinção, que é o núcleo do desenho:

\begin{description}
  \item[O valor do limiar] \textbf{muda} por litologia. É o objetivo da calibração.
  \item[A regra que produz o valor] \textbf{não muda}. É idêntica para todas as 45 litologias.
\end{description}

A regra adotada é a seguinte. Para cada sonda, constrói-se a curva que relaciona o limiar ao número de regiões marcadas, somando as três imagens de limiar da litologia. O limiar escolhido é o \textbf{joelho} dessa curva, definido como o ponto de máxima distância à corda que liga seus dois extremos, após normalização dos eixos \cite{kneedle_satopaa2011}.

Três aspectos da regra são fixados deliberadamente. A imagem de \textit{descoberta} \textbf{não} entra no cálculo: escolhida por ser a mais rica em feições, ela puxaria o limiar para cima, reintroduzindo o viés de seleção por visibilidade que o protocolo existe para evitar. O eixo do limiar é tomado em escala \textbf{logarítmica}, porque os escores concentram-se na década baixa --- medidos entre 0,011 e 0,770 em uma litologia de referência, com mediana entre 0,040 e 0,080 --- e, em escala linear, o joelho recairia sistematicamente sobre o primeiro ponto da curva. Por fim, a regra \textbf{não possui parâmetro livre}: não há constante a ser fixada por inspeção, o que elimina a possibilidade de ajustá-la para favorecer um resultado.

A regra é aplicada mecanicamente. Reconhece-se, porém, que uma curva pode enganar em casos particulares, e o autor pode divergir do valor produzido. Para que essa divergência não se torne invisível, a ferramenta de calibração registra \textbf{ambos} os valores --- o produzido pela regra e o adotado pelo autor --- junto de uma justificativa textual obrigatória. Os experimentos comparam sempre os valores da regra, de modo a preservar uma única variável manipulada; a taxa de concordância entre regra e autor é reportada como resultado à parte, e constitui evidência sobre a adequação da própria regra.

\section{Varredura de Limiares Fora de Linha}
\label{sec:varredura}

Calibrar exige comparar muitos limiares. Feito de modo ingênuo --- escolher um valor, executar o Professor, observar, ajustar e reexecutar --- cada iteração custa minutos de processamento em GPU, o que inviabiliza o protocolo da seção anterior em tempo humano.

A inspeção do código do arcabouço utilizado mostra que o limiar de confiança é aplicado como \textbf{filtro posterior} à inferência. O modelo produz máscaras e escores sem conhecimento do limiar; este apenas descarta as detecções cujo escore não o ultrapassa. A supressão de não-máximos ocorre depois do filtro, mas processa as detecções em ordem decrescente de escore e só remove uma detecção usando outra de escore \textbf{maior} --- de onde se conclui que incluir detecções de escore baixo não altera as decisões tomadas sobre as de escore alto.

Segue-se que executar o Professor uma única vez com o limiar no piso, armazenar escores e polígonos, e filtrar esse material posteriormente é \textbf{exatamente equivalente} a reexecutar o modelo em cada limiar --- não se trata de aproximação. O ciclo de calibração passa de minutos por iteração a milissegundos, e é o que torna viável exibir simultaneamente o efeito de um mesmo limiar sobre as quatro imagens da litologia.

A equivalência não foi assumida: ela é verificada experimentalmente, e o procedimento de verificação é descrito na Seção \ref{sec:verificacao} e seu resultado reportado no Capítulo 4.

\section{Referência de Avaliação}
\label{sec:referencia}

Toda métrica quantitativa de segmentação exige uma referência. Usar a saída do próprio Professor como referência mediria fidelidade de cópia, não acerto, e tornaria a comparação entre braços experimentais destituída de sentido --- cada braço teria um gabarito diferente. Adotam-se, portanto, duas fontes de referência complementares.

\subsection{Conjunto de referência anotado}

Constitui-se um conjunto de aproximadamente 50 imagens, retiradas da partição de \textbf{teste} e distribuídas entre litologias da faixa A. Essas imagens são anotadas manualmente pelo autor \textbf{antes} de qualquer inferência ser executada sobre elas. A ordem é parte da metodologia: anotar depois de ver a máscara do modelo ancora o anotador nela, e a anotação deixa de ser independente.

É preciso delimitar o que essa precaução garante e o que não garante. Ela garante que a anotação não é cópia da máscara do modelo. Ela \textbf{não} garante neutralidade: o autor desenvolveu o protocolo de calibração observando saídas do Professor durante meses, e essa exposição molda sua noção do que constitui uma anomalia. A independência é parcial, e o trabalho a declara como limitação em vez de apresentá-la como garantia de imparcialidade.

Sobre esse gabarito calculam-se IoU, precisão, \textit{recall} e taxa de falso positivo, definido operacionalmente como \textbf{região marcada pelo modelo sem qualquer sobreposição} --- IoU igual a zero --- com região alguma do conjunto de referência.

\subsection{Preferência pareada cega}

A segunda fonte é qualitativa e envolve um profissional do setor. A mesma chapa é apresentada com a marcação de dois braços experimentais, embaralhadas e sem identificação, e o profissional indica qual representa melhor o que trataria como defeito. O procedimento produz estatística de preferência, e não IoU.

Registra-se explicitamente que essa fonte é \textbf{complementar}: as hipóteses H1 e H2 são enunciadas contra o conjunto de referência anotado, e não contra o julgamento do profissional. Caso a preferência pareada seja realizada, ela reforça as conclusões; sua ausência não invalida as hipóteses.

\section{Experimento 1 --- Critério por Litologia \textit{versus} Critério Único}

O primeiro experimento avalia H1 e não envolve treinamento algum: compara duas configurações do Professor sobre as mesmas imagens, no mesmo modelo e no mesmo equipamento.

O ponto delicado do desenho é a definição do braço de controle. Uma configuração única escolhida arbitrariamente tornaria a comparação vazia, pois o resultado dependeria dessa escolha. O braço de controle é, portanto, \textbf{derivado por regra}, e não selecionado: aplica-se a \textbf{mesma regra} da Seção \ref{sec:regra}, alterando-se uma única coisa --- o escopo do material sobre o qual ela é calculada, conforme a Tabela \ref{tab:bracos_exp1}.

\begin{table}[ht]
  \caption{Braços do Experimento 1. A regra é a mesma; muda apenas seu escopo.}
  \label{tab:bracos_exp1}
  \centering
  \begin{tabular}{ll}
    \hline
    \bfseries Braço & \bfseries Escopo da curva que produz o limiar \\
    \hline
    Calibrado & As três imagens de limiar \textbf{daquela} litologia \\
    Único     & As três imagens de limiar de \textbf{todas} as litologias reunidas \\
    \hline
  \end{tabular}
  \legend{Fonte: elaborado pelo autor.}
\end{table}

Sonda por sonda, reúnem-se os escores de todas as litologias em um único conjunto, aplica-se a regra e obtém-se um limiar global por sonda. O conjunto de sondas do braço único é aquele que a mesma regra de descoberta admite sobre o material reunido.

Cabe justificar por que o braço único não é definido como a mediana dos limiares calibrados, alternativa aparentemente natural. Se o fosse, o resultado seria quase aritmético: o valor ótimo de cada litologia supera a mediana dos valores ótimos \textbf{nos dados dessa mesma litologia} por construção, o que caracteriza regressão à média e não achado experimental --- além de tornar o controle uma função do tratamento. A regra aplicada ao material reunido é o contrafactual honesto: \textit{e se todas as litologias tivessem sido tratadas como uma só?}, que é precisamente o que H1 pergunta. Esse desenho admite resultado negativo, e é isso que torna a hipótese falseável.

A avaliação de ambos os braços é feita contra o conjunto de referência da Seção \ref{sec:referencia}, reportando-se IoU, precisão, \textit{recall} e taxa de falso positivo, por braço e por litologia.

\section{Experimento 2 --- Especialistas \textit{versus} Generalista}

O segundo experimento avalia H2 e compara modelos Aluno treinados sobre as anotações do Professor, em três braços, descritos no Quadro \ref{quad:bracos_exp2}.

\begin{quadro}[ht]
  \caption{Braços do Experimento 2.}
  \label{quad:bracos_exp2}
  \centering
  \begin{tabularx}{\textwidth}{lX}
    \hline
    \bfseries Braço & \bfseries Descrição \\
    \hline
    Especialista & Um modelo por litologia, treinado apenas com as anotações daquela litologia \\
    Generalista  & Um único modelo, treinado com as anotações de todas as litologias \\
    Controle     & Um único modelo, treinado com anotações calibradas por litologia \\
    \hline
  \end{tabularx}
  \legend{Fonte: elaborado pelo autor.}
\end{quadro}

O braço de controle existe por necessidade de desenho: sem ele, a comparação entre especialistas e generalista alteraria simultaneamente duas variáveis --- a quantidade de modelos e a especificidade das anotações --- e nenhum resultado seria atribuível a uma delas.

Os resultados são reportados \textbf{por faixa} de volume, e não agregados, o que responde à segunda metade de H2.

Reconhece-se uma limitação de interpretação: cada modelo especialista é treinado com uma fração dos dados vistos pelo generalista, de modo que a comparação mede o efeito \textbf{líquido} da especialização em cada volume de dados, e não um efeito causal isolado do número de modelos. A estratificação por faixa é justamente o que torna esse efeito líquido interpretável.

\section{Ambiente Experimental e Reprodutibilidade}
\label{sec:verificacao}

Os experimentos são executados no ambiente descrito no Quadro \ref{quad:ambiente}.

\begin{quadro}[ht]
  \caption{Ambiente de execução dos experimentos.}
  \label{quad:ambiente}
  \centering
  \begin{tabularx}{\textwidth}{lX}
    \hline
    \bfseries Componente & \bfseries Versão \\
    \hline
    Linguagem            & Python 3.14 \\
    Arcabouço            & PyTorch 2.11 \cite{pytorch_paszke2019} com CUDA 12.8 \\
    Biblioteca de modelos& Ultralytics 8.4.61 \\
    Processador gráfico  & NVIDIA GeForce RTX 5060 Ti, 16 GB \\
    \hline
  \end{tabularx}
  \legend{Fonte: elaborado pelo autor.}
\end{quadro}

A equivalência da varredura fora de linha, afirmada na Seção \ref{sec:varredura}, é verificada por um procedimento automatizado: para cada sonda e para uma amostra de limiares, compara-se o resultado de executar o Professor com aquele limiar ao resultado de filtrar o material armazenado no mesmo valor. A comparação é feita sobre o número de detecções, sobre os escores e sobre os vértices dos polígonos, \textbf{sem tolerância numérica}. O procedimento é parte do repositório do trabalho e deve ser reexecutado sempre que a versão da biblioteca de modelos mudar, uma vez que a equivalência é propriedade da implementação e não do método.

Todos os artefatos de configuração --- o conjunto de sondas e limiares por litologia, os metadados de origem das imagens de calibração e o registro de calibração com a justificativa textual --- são versionados junto do código.

  \chapter[RESULTADOS]{RESULTADOS}

Este capítulo reporta os resultados obtidos até o momento. Ele começa pelas verificações que sustentam o método descrito no Capítulo 3 --- e que são pré-requisito para que os experimentos façam sentido --- e prossegue para os resultados dos experimentos propriamente ditos.

% =====================================================================
% TODO (estrutura): este capítulo está completo apenas nas seções 4.1 a 4.3.
% As seções do Experimento 1 e do Experimento 2 aguardam execução. Ver
% docs/roadmap.md no repositório para o estado corrente.
% =====================================================================

\section{Verificação da Equivalência da Varredura Fora de Linha}

A Seção \ref{sec:varredura} afirma que executar o Professor uma única vez com o limiar no piso e filtrar o material armazenado é exatamente equivalente a reexecutar o modelo em cada limiar. Como todo o protocolo de calibração depende dessa afirmação, ela foi verificada por execução, e não apenas por inspeção do código.

O procedimento fixou uma imagem e, para cada uma de quatro sondas, comparou duas grandezas em oito limiares distribuídos entre 0,005 e 0,500: o resultado de executar o Professor com aquele limiar e o resultado de filtrar, naquele mesmo valor, o material capturado com o limiar no piso de 0,001. A comparação abrangeu o número de detecções, os escores e os vértices dos polígonos, \textbf{sem tolerância numérica}.

As 32 comparações resultaram idênticas em todas as três grandezas. Não houve divergência em nenhum limiar, nem no número de regiões, nem em qualquer coordenada de polígono. A Tabela \ref{tab:d18} resume o número de detecções sobreviventes em cada configuração --- valor que, por ser idêntico nos dois métodos, é reportado em coluna única.

\begin{table}[ht]
  \caption{Detecções sobreviventes por limiar, idênticas entre execução real e varredura fora de linha.}
  \label{tab:d18}
  \centering
  \begin{tabular}{lrrrrrrrr}
    \hline
    \bfseries Sonda & \bfseries 0,005 & \bfseries 0,02 & \bfseries 0,05 & \bfseries 0,08 & \bfseries 0,12 & \bfseries 0,20 & \bfseries 0,35 & \bfseries 0,50 \\
    \hline
    \textit{crack}        & 168 & 165 & 118 & 76 & 41 & 20 & 9 & 4 \\
    \textit{vein}         & 164 & 149 & 71  & 34 & 21 & 11 & 7 & 4 \\
    \textit{Stain}        & 195 & 186 & 132 & 94 & 74 & 48 & 15 & 7 \\
    \textit{Dark patches} & 194 & 175 & 120 & 73 & 52 & 25 & 9 & 6 \\
    \hline
  \end{tabular}
  \legend{Fonte: elaborado pelo autor.}
\end{table}

O resultado tem duas consequências. A primeira é prática: o ciclo de calibração deixa de custar minutos de GPU por iteração e passa a custar milissegundos, o que é o que viabiliza submeter um mesmo limiar às quatro imagens da litologia simultaneamente. A segunda é metodológica: a varredura fora de linha não introduz aproximação alguma, de modo que nenhum resultado obtido por ela precisa ser qualificado como estimativa.

Ressalva necessária: a equivalência é propriedade da \textbf{implementação} verificada, não do método em abstrato. Uma versão futura da biblioteca que aplicasse o limiar antes da supressão de não-máximos, ou que o incorporasse ao próprio modelo, invalidaria a conclusão. Por isso o procedimento de verificação integra o repositório e deve ser reexecutado a cada atualização de versão.

\section{Caracterização dos Escores do Professor}

A escolha da escala logarítmica na regra de limiar (Seção \ref{sec:regra}) decorre da distribuição observada dos escores, medida sobre a imagem de descoberta de uma litologia da faixa A e reportada na Tabela \ref{tab:escores}.

\begin{table}[ht]
  \caption{Distribuição dos escores do Professor por sonda, com o limiar no piso.}
  \label{tab:escores}
  \centering
  \begin{tabular}{lrrrr}
    \hline
    \bfseries Sonda & \bfseries Detecções & \bfseries Mínimo & \bfseries Mediana & \bfseries Máximo \\
    \hline
    \textit{crack}        & 168 & 0,016 & 0,074 & 0,699 \\
    \textit{vein}         & 164 & 0,013 & 0,045 & 0,694 \\
    \textit{Stain}        & 195 & 0,012 & 0,076 & 0,683 \\
    \textit{Dark patches} & 194 & 0,011 & 0,061 & 0,766 \\
    \hline
  \end{tabular}
  \legend{Fonte: elaborado pelo autor.}
\end{table}

Os escores percorrem quase duas ordens de grandeza, de 0,011 a 0,766, mas a mediana situa-se entre 0,045 e 0,076 em todas as sondas: a massa da distribuição concentra-se na década baixa e a cauda superior é rala. Em uma grade linear de limiares, portanto, a quase totalidade dos pontos cairia em região sem detecção alguma, e o joelho da curva recairia sistematicamente sobre seu primeiro ponto --- daí a normalização logarítmica do eixo.

O dado também recomenda cautela na transposição de valores da literatura de pseudo-rotulagem. Os limiares elevados ali reportados referem-se a escores de detectores supervisionados calibrados, cuja escala não corresponde à observada aqui. O que se transporta daquela literatura é a direção do compromisso --- privilegiar precisão --- e não a magnitude do corte.

\section{Geometria dos Polígonos Produzidos pelo Professor}

A conversão das máscaras do Professor em polígonos no formato de anotação apresenta uma particularidade com consequência sobre a qualidade do rótulo. Quando a máscara de uma detecção possui mais de um contorno desconexo, a implementação utilizada funde todos eles em uma única poligonal, ligando-os por trajetos de ida e volta.

A medição sobre a imagem de descoberta de uma litologia da faixa A, com a sonda \textit{crack} no limiar de 0,08, encontrou 45 detecções com mais de um contorno em um total de 76. Como os trajetos de ligação têm área praticamente nula, o efeito agregado é pequeno: a área total dos polígonos supera a das máscaras em 2,4\%, e a sobreposição entre a união das máscaras e a união dos polígonos atinge IoU de 0,928. O caso individual, porém, não é igualmente benigno: em uma detecção composta por 18 contornos, a área do polígono alcançou 2,5 vezes a da máscara correspondente.

A implicação é que o efeito é desprezível quando se avalia a cobertura agregada, mas não quando uma detecção fragmentada é tomada isoladamente como instância de treinamento. O tratamento dessa geometria --- seja pela retenção do maior contorno, seja pela separação dos contornos em instâncias distintas --- integra o pós-processamento do Professor e é registrado como etapa metodológica a definir, e não como limitação incidental.

\section{Experimento 1 --- Critério por Litologia \textit{versus} Critério Único}

% TODO: executar e redigir. Depende da calibração das litologias da faixa A e da
% anotação do conjunto de referência (Seção \ref{sec:referencia}).
% Reportar, por braço e por litologia: IoU, precisão, recall e taxa de falso
% positivo; incluir o par de figuras "critério único x critério calibrado" sobre
% a mesma chapa; e reportar a taxa de concordância entre a regra e o autor.

\section{Experimento 2 --- Especialistas \textit{versus} Generalista}

% TODO: executar e redigir, faixa por faixa. Depende da geração do conjunto de
% treinamento pelo Professor, do pós-processamento discutido na Seção 4.3 e do
% treinamento dos três braços descritos no Quadro \ref{quad:bracos_exp2}.

  \chapter[CONSIDERAÇÕES FINAIS]{CONSIDERAÇÕES FINAIS}

\section{Síntese}

Este trabalho partiu de uma constatação sobre a inspeção de chapas de rochas ornamentais: o critério que separa uma feição natural de um defeito comercial não existe por escrito. Ele é real, é aplicado milhares de vezes por dia e determina o preço das peças --- mas vive implícito na experiência de cada operador, onde não pode ser consultado, comparado entre turnos nem auditado por um cliente que conteste uma classificação.

A proposta do ARIA não foi eliminar a arbitrariedade dessa fronteira, o que seria promessa que nenhum sistema pode cumprir, dado que a fronteira depende de convenção comercial e não de propriedade física. A proposta foi \textbf{deslocá-la}: retirá-la da cabeça do inspetor e materializá-la em um conjunto de sondas textuais e limiares de confiança por litologia, gravado em arquivo versionável, produzido por uma regra idêntica para todas as litologias e acompanhado do registro de cada divergência entre a regra e o julgamento do autor.

% TODO: acrescentar aqui, quando os experimentos rodarem, dois a três parágrafos
% sintetizando os resultados de H1 e H2 -- inclusive se o resultado for negativo.
% Resultado negativo em H1 (critério calibrado equivalente ao critério único) é
% achado legítimo e deve ser reportado como tal, não omitido: o desenho descrito
% na Seção 3.8 foi construído justamente para admitir essa possibilidade.

Independentemente do desfecho dos experimentos, dois resultados já se firmaram. O primeiro é a demonstração de que o limiar de confiança do Professor é um filtro aplicado após a inferência, o que torna a varredura de limiares fora de linha exatamente equivalente a reexecutar o modelo --- propriedade verificada experimentalmente, sem tolerância numérica, e que reduz o ciclo de calibração de minutos para milissegundos. O segundo é o próprio protocolo de calibração: a separação entre a imagem que descobre \textit{quais} sondas entram e as três que submetem um \textit{único} limiar a casos sutil, típico e forte, com a distinção explícita entre o valor do limiar, que varia por litologia, e a regra que o produz, que não varia.

\section{Limitações}

O trabalho reconhece as seguintes limitações.

\begin{description}

  \item[Independência parcial da referência] O conjunto de referência foi anotado antes de qualquer inferência sobre as mesmas imagens, o que impede que a anotação seja cópia da máscara do modelo. Isso não equivale a neutralidade: o autor desenvolveu o protocolo observando saídas do Professor, e essa exposição molda sua noção de anomalia. A referência é independente da predição, não do processo que a gerou.

  \item[As sondas não identificam o tipo da anomalia] O sistema marca regiões que responderam a determinada chave lexical. Ele não afirma, e este trabalho não valida, que a região marcada por \textit{crack} seja uma fissura. Toda a rotulagem é binária, e qualquer leitura semântica das sondas extrapola o que foi verificado.

  \item[Efeito líquido, não isolado, da especialização] No Experimento 2, cada modelo especialista é treinado com uma fração dos dados vistos pelo generalista. A comparação mede o efeito conjunto da especialização e do volume de dados por modelo, e é a estratificação por faixa que torna esse efeito interpretável --- não um controle que o isole.

  \item[A equivalência da varredura é propriedade da implementação] A verificação reportada vale para a versão da biblioteca utilizada. Alterações na ordem entre filtragem e supressão de não-máximos invalidariam a conclusão, razão pela qual o procedimento de verificação foi mantido no repositório.

  \item[Integração entre estágios não demonstrada de ponta a ponta] O classificador de litologias é modelo já treinado e avaliado em trabalho anterior do grupo \cite{deepstoneai_moreira2025}, mas o roteamento automático que o conecta ao segundo estágio não é demonstrado nesta versão. O desempenho em quadros por segundo do estágio Aluno, igualmente, não é medido --- e por isso nenhuma afirmação de operação em tempo real é feita com base em medição própria.

  \item[Cobertura parcial das faixas de volume] O desenho do Experimento 2 prevê execução faixa a faixa. As faixas não alcançadas dentro do prazo do trabalho são declaradas explicitamente, e não silenciadas por agregação dos resultados.

\end{description}

\section{Trabalhos Futuros}

\begin{description}

  \item[Rotulagem multiclasse] Preservar a distinção entre sondas no treinamento do Aluno exigiria validar que cada sonda de fato identifica o fenômeno que seu nome sugere --- validação que este trabalho deliberadamente não faz. Como os identificadores por sonda já são gravados nos arquivos de anotação, a extensão não exigiria reprocessar o conjunto de dados.

  \item[Aprendizado ativo] O arranjo Professor--Aluno admite um ciclo em que o Professor reexamina, fora de linha, as predições de baixa confiança do Aluno em produção, realimentando o conjunto de treinamento. O ciclo não integra esta versão e não é reivindicado como contribuição.

  \item[Faixas de volume remanescentes] A execução das faixas não alcançadas completa a curva que relaciona o ganho da especialização ao volume de dados por litologia, que é a resposta à segunda metade de H2.

  \item[Sensibilidade a \textit{prompts} contextuais] O trabalho que originou as representações multimodais empregadas reporta que \textit{prompts} contextuais superam palavras isoladas em tarefas \textit{zero-shot} \cite{clip_radford2021}. Como as sondas aqui adotadas são majoritariamente palavras isoladas, e algumas são lexicalmente ambíguas, uma análise de sensibilidade comparando as duas formulações é extensão natural e de baixo custo --- bastaria alterar o arquivo de configuração.

  \item[Integração de ponta a ponta e medição de latência] Implementar o roteamento automático entre os três estágios e medir a latência real do sistema completo, substituindo afirmações qualitativas de velocidade por números.

\end{description}

    
  % --------------------------------------------------- %
  % ELEMENTOS POS-TEXTUAIS %
  % --------------------------------------------------- %
  \postextual
  
  % Referencias Bibliográficas
  \bibliographystyle{abntex2-alf} % Autor-Data
  \bibliography{bibliografia}
    
  % Apêndices
  \begin{apendicesenv}
    %%%\chapter{{\normalfont Explicação}}
\label{apendice: a}

É um elemento opcional. É um documento elaborado pelo próprio autor com o objetivo de completar sua argumentação, sem que haja prejuízo para a unidade do trabalho (ASSOCIAÇÃO BRASILEIRA DE NORMAS TÉCNICAS, 2011b. Deve ser precedido da palavra APÊNDICE (ex:APÊNDICE A, APÊNDICE B), identificado por letras maiúsculas consecutivas, travessão e pelo respectivo título. 
  \end{apendicesenv}

  % Anexos
  \begin{anexosenv}
    %%%\chapter{{\normalfont Explicação}}
\label{anexo: a}

É um elemento opcional. Não é elaborado pelo próprio autor mas sim por outras pessoas e constitui-se de suportes elucidativos e ilustrativos importantes para a compreensão do texto (ASSOCIAÇÃO BRASILEIRA DE NORMAS TÉCNICAS, 2011b). Havendo mais de um anexo, sua identificação deve ser feita por letra maiúscula ou algarismo arábico (ex:ANEXO A, ANEXO B), identificado por letras maiúsculas consecutivas,travessão e pelo respectivo título. 
    %%%\chapter{{\normalfont EasyReview}}

\lstset{
		keywordstyle=\bfseries\color{purple!60!black},
		commentstyle=\itshape\color{green!40!black},
% 		identifierstyle=\itshape,
		stringstyle=\itshape,
		showstringspaces=false,
		escapeinside={\%*}{*)},
		columns=flexible,
		breaklines=true,
breakindent=0pt}

\renewcommand\lstlistingname{Example}
\newcommand{\keyword}[1]{\textbf{\color{purple!60!black}{#1}}}

% ------
\textbf{\uppercase {The alert command}}

% ------
Command intended to claim author's attention to a given part of the text. In the following, it is possible to see an example:

\begin{center}
\begin{minipage}[ht]{0.45\textwidth}
\begin{lstlisting}[language=tex]
A text without the alert command. %*\keyword{$\backslash$alert}*){A text with the alert command}.
\end{lstlisting}
\end{minipage}
\hspace{10pt}
\begin{minipage}[ht]{0.45\textwidth}
A text without the alert command. \alert{A text with the alert command}.
\end{minipage}
\end{center}

% ------
\textbf{\uppercase {The highlight command}}

% ------
Command intended to claim author's attention to a given part of the text in a different to the ``alert'' command. In the following, it is possible to see an example:

\begin{center}
\begin{minipage}[ht]{0.45\textwidth}
\begin{lstlisting}[language=tex]
A text without the highlight command. %*\keyword{$\backslash$highlight}*){A text with the highlight command}.
\end{lstlisting}
\end{minipage}
\hspace{10pt}
\begin{minipage}[ht]{0.45\textwidth}
A text without the highlight command. \highlight{A text with the highlight command}.
\end{minipage}
\end{center}

% ------
\textbf{\uppercase {The remove command}}

% ------
Command which an author suggest to remove a given part of the text. In the following, it is possible to see an example:

\begin{center}
\begin{minipage}[ht]{0.45\textwidth}
\begin{lstlisting}[language=tex]
This text is not to be removed. %*\keyword{$\backslash$remove}*){This text is to be removed}.
\end{lstlisting}
\end{minipage}
\hspace{10pt}
\begin{minipage}[ht]{0.45\textwidth}
This text is not to be removed. \remove{This text is to be removed}.
\end{minipage}
\end{center}

% ------
\textbf{\uppercase {The add command}}

% ------
Command which an author suggest to add new text in a given part of the text. In the following, it is possible to see an example:

\begin{center}
\begin{minipage}[ht]{0.45\textwidth}
\begin{lstlisting}[language=tex]
This text was already in the text. %*\keyword{$\backslash$add}*){This text is been added now}.
\end{lstlisting}
\end{minipage}
\hspace{10pt}
\begin{minipage}[ht]{0.45\textwidth}
This text was already in the text. \add{This text is been added now}.
\end{minipage}
\end{center}

% ------
\textbf{\uppercase {The replace/substitute command}}

% ------
Both commands are equivalent. It shall be used when an author suggest to replace a given part of the text for a newer one. In the following, it is possible to see an example:

\begin{center}
\begin{minipage}[ht]{0.45\textwidth}
\begin{lstlisting}[language=tex]
%*\keyword{$\backslash$replace}*){This part of the text needs to be replaced}{for this newer part of the text}.
\end{lstlisting}
\end{minipage}
\hspace{10pt}
\begin{minipage}[ht]{0.45\textwidth}
\replace{This part of the text needs to be replaced}{for this newer part of the text}.
\end{minipage}
\end{center}

% ------
\textbf{\uppercase {The commentreview command}}

% ------
Command intended to claim author's attention to a given part of the text, giving some comments to provide more information. In the following, it is possible to see an example:

\begin{center}
\begin{minipage}[ht]{0.45\textwidth}
\begin{lstlisting}[language=tex]
%*\keyword{$\backslash$commentreview}*){This text will receive a comment.}{This is the comment I have!}.
\end{lstlisting}
\end{minipage}
\hspace{10pt}
\begin{minipage}[ht]{0.45\textwidth}
\commentreview{This text will receive a comment.}{This is the comment I have!}.
\end{minipage}
\end{center}
  \end{anexosenv}

  % Índice Remissivo
  \phantompart
  \printindex

\end{document}